\documentclass[11pt]{article}
\usepackage[margin=1.15in]{geometry}
\usepackage{amsmath,amssymb,url}
\pagestyle{empty}
\begin{document}

\noindent
Johan Hoffman and Claes Johnson\\
KTH Royal Institute of Technology\\
SE-100 44 Stockholm, Sweden\\[1.2em]
\today\\[1.6em]
%
\noindent
Clay Mathematics Institute\\
Scientific Advisory Board\\[1.8em]

\noindent Dear Members of the Board,\\[0.8em]

We write to present a solution to the Navier--Stokes Millennium Problem, in a form which differs from
the one the official statement anticipates, and to explain why we believe the difference is necessary
rather than evasive. We should say at once that we regard the official formulation as sound
mathematics and as having been extraordinarily productive; what we argue is that it is incomplete, in
a way that twenty-five years of subsequent work has made visible and that could not have been seen
when it was written. We enclose the article \emph{Reformulation and Solution of the Clay
Navier--Stokes Problem}.

\paragraph{What we claim.} Two things, offered together and standing or falling separately.

First, a \textbf{reformulation}. We ask for \emph{constructive existence and qualitative uniqueness of
weak solutions to the incompressible Navier--Stokes equations at high Reynolds number}, for given data:
a domain $\Omega\subset\mathbb{R}^3$, initial data, a forcing, and boundary conditions. By
\emph{constructive} existence we mean existence established by exhibition --- a solution produced, with
the mesh, the discretisation, the measured residual, the computed stability factors and, for each
output, a value with a bound on its error. By \emph{qualitative} uniqueness we mean that the data
determine the flow structure and the mean-value outputs, with the line between what is determined and
what is not drawn quantitatively, functional by functional, rather than asserted.

Second, a \textbf{solution of that problem}: computed weak solutions for specific data, with the error
in each output bounded by duality against a computed solution of the adjoint equations, and with mean
drag and lift converging to a few per cent on meshes on which the velocity field does not converge and
never will. The computational record is in the literature and is cited in the article.

\paragraph{Why not the official formulation.} The official statement offers four alternatives, and asks
whether smooth data give a solution that stays smooth or whether the velocity blows up. The flows of
physical interest are neither. They are turbulent: not smooth, since velocity gradients grow like
$\nu^{-1/2}$ and are enormous at any Reynolds number of interest, and not singular, since the velocity
stays bounded. An alternative which omits the case that occurs can be settled without
bearing on fluids.

It is worth putting the point at its strongest, since it is not a statement about which solutions are
observed. A high Reynolds number flow set going laminar develops shear instabilities, and these steepen
the velocity gradients they feed on until the gradients are sharp enough for viscosity to act ---
order $\nu^{-1/2}$, on structures of order $\nu^{3/4}$. The laminar state therefore has no
continuation: it cannot persist, because the instability grows; it cannot escape by blowing up, because
the steepening is arrested as soon as viscosity bites; and it cannot dissipate in the laminar state what
the large scales supply to it. One route remains, and the flow takes it. It becomes turbulent in order
to continue to exist. Turbulence is thus not a pathology of the equations but the mode in which their
solutions exist at high Reynolds number, and it is the case the official formulation treats as
failure.

We would add an observation about mathematical practice. A theorem and its proof are constructed in
parallel; the statement is adjusted as the argument develops until the two fit. A formulation fixed in
advance of any proof, by a third party, and left standing for twenty-five years is therefore an unusual
object, and one should not be surprised if what eventually satisfies it is an argument built to fit
rather than a discovery about fluids.

\paragraph{The announced solution, and why it does not close the matter.} We take the September 2026
solution announced by OpenAI to be a serious piece of work and we do not dispute it. We observe two
things about it.

Its Theorem~1.1 constructs, for every $\nu>0$, a force $f\in C^\infty_c$, compactly supported in space
and time, and fields starting from rest whose kinetic energy stays bounded while the velocity becomes
unbounded at $t=1$. The order of construction is what we would draw attention to: the velocity and
pressure are built first, to have the required growth, and the force is then defined as what remains
--- the paper's final section sets out to give the local field ``zero initial datum and compact spatial
support, while making its momentum residual the restriction of a force in
$C^\infty_c$''. The force is thus the residual of a flow chosen in advance, arranged to be smooth,
rather than a physical forcing mechanism. And the statement it proves is (C), as the paper itself says --- ``this establishes alternative (C) in
the Millennium problem statement for Navier--Stokes as stated by Fefferman'', with (D) following from
compact support. (C) is not the negation of (A): (A) and (B) are posed
with $f\equiv0$ and quantified over all admissible data, while (C) and (D) ask for the existence of data
\emph{together with an admissible forcing} for which smoothness fails. A blowup driven by an external
force shows that something can be done \emph{to} a fluid; (A) asks whether a fluid does it to itself.
So the prize may be won without the question everyone means by the Navier--Stokes problem being
settled.

There is also a quantitative point. Navier--Stokes rescales, so one unit-viscosity blowup generates the
whole family, and the exponents follow from which normalisation is imposed. Held to a fixed admissible
forcing class --- which is the case that bears on the official statement, since it is that class the
formulation fixes --- the singular solution has amplitude $\nu^{1/3}$ on support $\nu^{2/3}$, so the
family tends to $u\equiv0$ and the singularity vanishes in the limit where fluids become turbulent.
Held instead at fixed amplitude, the force diverges like $\nu^{-1}$ and leaves the class. A singularity
present at every fixed viscosity and absent as $\nu\to0$ is a property of the admissible class rather
than of the flow.

\paragraph{On the condition of prior publication.} The resolution we present is not new and is not
offered to the Board for the first time. It is documented in the series of articles and the book cited
below, in refereed journals of high standard, starting in 2007. The requirement of publication in a
refereed journal, with two years elapsed since, is therefore met, and met by a margin approaching two
decades.

The remaining condition, that the work enjoy general acceptance in the community, is not ours to
declare.

One thing should be plain. What the cited work solves is the reformulated problem. Whether it thereby
answers the official one is the question the reformulation raises, and it is a question about the
official formulation rather than about the solution.

We would be glad to supply computational detail, or to discuss the formulation with anyone the Board
cares to nominate.\\[1.4em]

\noindent Yours faithfully,\\[2.2em]

\noindent Johan Hoffman \hfill Claes Johnson\\
\noindent KTH Royal Institute of Technology \hfill KTH Royal Institute of Technology

\vspace{2em}
\noindent\rule{\textwidth}{0.4pt}

\paragraph{The published record.}
\begin{enumerate}\itemsep2pt
\item J.~Hoffman and C.~Johnson, \emph{Computational Turbulent Incompressible Flow}, Applied
  Mathematics: Body and Soul, vol.~4, Springer, 2007.
\item J.~Hoffman, \emph{Efficient computation of mean drag for the subcritical flow past a circular
  cylinder using general Galerkin G2}, Int.\ J.\ Numer.\ Meth.\ Fluids \textbf{59} (2009) 1241--1258.
\item J.~Hoffman and C.~Johnson, \emph{Resolution of d'Alembert's paradox}, J.\ Math.\ Fluid Mech.\
  \textbf{12} (2010) 321--334.
\item J.~Hoffman, J.~Jansson and R.~Vilela de Abreu, \emph{Adaptive modeling of turbulent flow with
  residual based turbulent kinetic energy dissipation}, Comput.\ Methods Appl.\ Mech.\ Engrg.\
  \textbf{200} (2011) 2758--2767.
\item J.~Hoffman, J.~Jansson, N.~Jansson and R.~Vilela de Abreu, \emph{Towards a parameter-free method
  for high Reynolds number turbulent flow simulation based on adaptive finite element approximation},
  Comput.\ Methods Appl.\ Mech.\ Engrg.\ \textbf{288} (2015) 60--74.
\item J.~Hoffman, J.~Jansson and C.~Johnson, \emph{New theory of flight}, J.\ Math.\ Fluid Mech.\
  \textbf{18} (2016) 219--241.
\end{enumerate}

\end{document}
